\documentclass[10pt,a4paper]{article}
\usepackage[utf8]{inputenc}
\usepackage{gdmath}
\usepackage{amscd}

\usepackage[margin=2.4cm]{geometry}

\title{Some Proofs in Category Theory: The Yoneda Lemma}
\author{Georgy Dunaev \\ \texttt{georgedunaev@gmail.com}}
\date{}

\begin{document}
\sloppy
\maketitle

\begin{abstract}
This monograph provides a self-contained, formally rigorous development of the Yoneda Lemma and its foundational consequences, starting from first-order predicate calculus with equality and axiomatic set theory. In Section 1, we present a Hilbert-style axiomatization of classical predicate logic with equality, prove the Deduction Metatheorem and quantifier distribution laws, and develop set-theoretic foundations including class and set intersections (\interclass K and \inter K), uniqueness and dual quantifiers (!x, \exists !x, ?x), the definite description binder \The x. P(x), Kuratowski ordered pairs, total and partial function application, dependent sums (\Sigma R) and dependent products (\Pi R) with \beta/\eta-conversion, functional extensionality, and transitive ordinals (\Tr and \TOrd). In Section 2, we build category theory directly upon these set-theoretic foundations: we verify the category axioms for \Set, opposite categories C^\op, and functor categories [C, D]; prove both the contravariant and covariant Yoneda Lemmas (\Nat(h_A, G) \cong G(A) and \Nat(h^A, F) \cong F(A)), their naturality in both arguments, and the full faithfulness of the Yoneda embedding; and develop representable functors, universal elements, the category of elements \int_C G, Grothendieck's terminality criterion, limits and colimits, continuity of Hom-functors, the Density Theorem (G \cong \mathrm{colim}_{\int_C G} h_A), adjoint functors (L \dashv R, RAPL/LAPC), Cartesian closure of presheaves, and Kan extensions.
\end{abstract}

\tableofcontents

\section{Logic and Set Theory}

%% EQUIVALENCE OF DEPENDENT PRODUCT DEFINITIONS %%
\begin{theorem}
For $f \in \Pow(\Sigma R)$, the compact condition $\Fun(f) \land \dom(R) \subseteq \dom(f)$ is equivalent to $\Rel(f) \land \Fun(f) \land \dom(f) = \dom(R) \land \forall x \in \dom(R).\ f(x) \in R(x)$.
\end{theorem}
\begin{proof}
[Proof content will be added here]
\end{proof}
%% SET-HOOD OF ΣR AND ΠR %%
\begin{theorem}
$\Sigma R \in \Set$ and $\Pi R \in \Set$ whenever $\dom(R) \in \Set$ and $\forall x \in \dom(R).\ R(x) \in \Set$.
\end{theorem}
\begin{proof}
[Proof content will be added here]
\end{proof}
%% β/η-LAWS %%
\begin{theorem}
For set-theoretic $\lambda$-abstraction, $\beta$- and $\eta$-conversion hold.
\end{theorem}
\begin{proof}
[Proof content will be added here]
\end{proof}\label{sec:logic_set}

\subsection{Foundations and Size Conventions}\label{sec:foundations}
To treat both small mathematical structures and large categorical constructions (such as the category of all sets \Set, opposite categories C^\op, functor categories [C, D], and categories of elements \int_C G) without set-theoretic paradoxes, we work in \textbf{von Neumann--Bernays--G"odel (\mathrm{NBG}) set theory} with the following size conventions:

1. All sets are small (belong to the universe \Set).
2. Classes are large collections that may not be sets.
3. The universe of sets \Set is a proper class.
4. We use the von Neumann hierarchy of sets to avoid paradoxes.
5. The axiom of choice holds for sets but not necessarily for proper classes.

We write $x \in y$ for membership, $x \subseteq y$ for subset, and $x \in y$ for proper subset. The power set of a set $x$ is denoted $\Pow(x) = \{y: y \subseteq x\}$. The empty set is $\emptyset = \{ \}$. We use the von Neumann ordinals $\Ord$ with $0 = \emptyset$, $1 = \{0\}$, $2 = \{0,1\}$, etc., and the class of transitive sets $\Tr = \{x: x \subseteq \Pow(x)\}$. The class of transitive ordinals is $\TOrd = \Ord \cap \Tr$.